\chapter{Sharing and Division}\label{ch:g3:division}

Twenty marbles shared fairly among four children: how many each?
Division answers all the sharing questions --- and the grouping
questions too --- and it is multiplication's twin: one undoes the
other.

\section{Two faces of division}

\begin{definition}[Division]\label{def:g3:division:def}
The division $20 \div 4$ answers two kinds of questions:
\begin{itemize}
  \item \emph{sharing}: $20$ marbles shared among $4$ children ---
        how many marbles \emph{each}? ($5$ each);
  \item \emph{grouping}: $20$ marbles put in bags of $4$ --- how many
        \emph{bags}? ($5$ bags).
\end{itemize}
Both questions have the same answer, because both undo the same
multiplication: $4 \times 5 = 20$.
\end{definition}

\begin{omfigure}
\begin{tikzpicture}[scale=0.55]
  \foreach \g in {0,1,2,3} {
    \draw[omExo, rounded corners] (\g*4.9-0.65,-0.7) rectangle
      (\g*4.9+3.85,0.7);
    \foreach \c in {0,...,4} \fill[omDef] (\g*4.9+\c*0.8,0) circle (0.27);
  }
  \node[below, font=\small] at (8.95,-1.15)
    {$20$ marbles, $4$ equal groups: $5$ in each group};
\end{tikzpicture}

{\small Sharing $20$ into $4$ equal groups. Grouping by $4$ instead
would give $5$ groups --- the same division.}
\end{omfigure}

\begin{method}[Dividing with the tables]\label{met:g3:division:tables}
To compute $42 \div 6$, ask: ``$6$ times \emph{what} makes $42$?''
The table of $6$ answers: $6 \times 7 = 42$, so $42 \div 6 = 7$.
Every division question is a multiplication question read
backwards.
\end{method}

\section{When it does not come out even}

\begin{definition}[Remainder]\label{def:g3:division:remainder}
Sharing $23$ marbles among $4$ children gives $5$ marbles each ---
and $3$ marbles \emph{left over}, not enough for another round. We
write
\[
23 = 4 \times 5 + 3 ,
\]
and say: the \emph{quotient} is $5$, the
\emph{remainder}\index{remainder} is $3$. The remainder is always
smaller than the number of children --- otherwise another round of
sharing would be possible.
\end{definition}

\begin{omfigure}
\begin{tikzpicture}[scale=0.55]
  \foreach \g in {0,1,2,3} {
    \draw[omExo, rounded corners] (\g*4.9-0.65,-0.7) rectangle
      (\g*4.9+3.85,0.7);
    \foreach \c in {0,...,4} \fill[omDef] (\g*4.9+\c*0.8,0) circle (0.27);
  }
  \foreach \c in {0,1,2} \fill[omThm] (19.7+\c*0.8,0) circle (0.27);
  \node[above, font=\footnotesize, omThm] at (20.5,0.5) {left over};
  \node[below, font=\small] at (10.2,-1.15)
    {$23 = 4 \times 5 + 3$: quotient $5$, remainder $3$ (in red)};
\end{tikzpicture}

{\small Sharing $23$ among $4$: after $5$ rounds, $3$ marbles remain
--- too few for a $6$th round.}
\end{omfigure}

\begin{example}\label{ex:g3:division:remainder}
Divide $58$ by $7$. Scan the table of $7$:
$7 \times 8 = 56$ is the biggest product that fits under $58$
($7 \times 9 = 63$ is too much). So
\[
58 = 7 \times 8 + 2 :
\]
quotient $8$, remainder $2$. Check: $56 + 2 = 58$, and $2 < 7$.
\end{example}

\begin{method}[Interpreting the remainder]\label{met:g3:division:interpret}
After dividing, return to the story:
\begin{enumerate}
  \item ``how many each / how many full groups?'' --- the quotient;
  \item ``how many left over?'' --- the remainder;
  \item ``how many boxes to hold everything?'' --- the quotient
        \emph{plus one} if the remainder is not zero.
\end{enumerate}
\end{method}

\begin{example}\label{ex:g3:division:interpret}
$58$ students go canoeing, $7$ per canoe. $58 = 7 \times 8 + 2$:
eight canoes are full, and $2$ students remain --- they need a canoe
too. So $9$ canoes must be booked.
\end{example}

\section{Halves, thirds, quarters}

\begin{definition}[Half, third, quarter]\label{def:g3:division:half}
The \emph{half} of a number is the result of dividing it by $2$; the
\emph{third}, by $3$; the \emph{quarter}, by $4$. For instance the
half of $18$ is $9$, the third of $18$ is $6$, the quarter of $36$
is $9$.
\end{definition}

\begin{example}[Doubles and halves go together]\label{ex:g3:division:double}
The half of $46$: since $46 = 40 + 6$, half of each part gives
$20 + 3 = 23$. Check with the double: $23 + 23 = 46$. Knowing
doubles by heart ($25 \to 50$, $45 \to 90$, $150 \to 300$) makes
halves instantaneous.
\end{example}

\section{Exercises}

\begin{exercise}[$\star$]\label{exo:g3:division:1}
Compute using the tables backwards: $36 \div 4$; \quad $63 \div 9$;
\quad $40 \div 5$; \quad $56 \div 8$.
\end{exercise}

\begin{exercise}[$\star$]\label{exo:g3:division:2}
For each division, write the equality $a = b \times q + r$ and give
the quotient and the remainder: $17 \div 5$; \quad $30 \div 4$; \quad
$50 \div 6$.
\end{exercise}

\begin{exercise}[$\star$]\label{exo:g3:division:3}
$28$ marbles are shared among $6$ children. How many marbles each,
and how many left over? Draw the sharing if it helps.
\end{exercise}

\begin{exercise}[$\star$]\label{exo:g3:division:4}
Compute: the half of $68$; the third of $27$; the quarter of $48$;
the half of $150$.
\end{exercise}

\begin{exercise}[$\star$]\label{exo:g3:division:5}
True or false? Check with a multiplication.
``$45 \div 5 = 9$''; \quad ``$52 \div 6 = 8$ remainder $4$'';
\quad ``$70 \div 8 = 9$ remainder $2$''.
\end{exercise}

\begin{exercise}[$\star$]\label{exo:g3:division:6}
$32$ photos are glued into an album, $4$ per page. How many pages are
used? Which face of division is this (sharing or grouping)?
\end{exercise}

\begin{exercise}[$\star$]\label{exo:g3:division:7}
A rope of $75$ cm is cut into pieces of $9$ cm. How many full pieces,
and how long is the leftover bit?
\end{exercise}

\begin{exercise}[$\star$]\label{exo:g3:division:8}
$50$ children go on a trip in cars of $4$ seats. How many cars are
needed? (Careful: everyone must ride!)
\end{exercise}

\begin{exercise}[$\star$]\label{exo:g3:division:9}
Find the missing numbers: $? \div 7 = 6$; \quad
$54 \div ? = 6$; \quad the half of $?$ is $35$.
\end{exercise}

\begin{exercise}[$\star\star$]\label{exo:g3:division:10}
Amina shares her $47$ stickers among her $5$ friends, as fairly as
possible, and keeps the leftover for herself. How many stickers does
each friend get, and how many does Amina keep?
\end{exercise}

\begin{exercise}[$\star\star$]\label{exo:g3:division:11}
When a number is divided by $6$, what are all the possible
remainders? What is the biggest number whose division by $6$ has
quotient $7$?
\end{exercise}
